\section{讨论与结论}\label{sec:discussion}

窗口重用改变的是派生特征的联合分布，而不只是使有效样本量下降。即使原始观测独立同分布，重叠窗口仍会共享观测，派生特征因而可能相关；全局置换保留经验边际分布，却打散原顺序中的滞后协方差。命题~\ref{prop:counterexample}给出这种层级错配的有限样本反例，$3.974$ 只刻画标量窗口均值的方差比，不能外推为高维最大扫描的误报率倍数。因此，科学零假设成立不能自动提供窗口层面的排列不变性。

“降低误报”不等于校准。M1 在全部 56 个匹配组合中降低拒绝率，却在非重叠端系统保守、在强重叠端仍残留反保守；把方向相反的局部偏差汇总，可能得到接近 $0.05$ 的均值而掩盖严重失真，故逐格拒绝率与 Wilson 分类才是主要证据。表示敏感性与误报控制是两个维度：功效反转说明固定 BPI 保留了矩摘要未捕获的几何信息，但这既不能替代校准评价，也不能把全部检出优势归因于一维同调通道。

本研究的局限在于：结论限于固定 $W$、两个合成数据生成过程、固定 BPI 与 M1 配置、单变点最大扫描；未纳入自然数据、更多数据生成过程、多变点递归与表示超参数遍历；本文没有给出 BPI 扫描的通用有限样本定理，逐格分类为点态判据，M1 的带宽规则也未针对每个步长优化。对这些条件的推广需要新的理论与独立确认实验。
